\documentclass[11pt]{article}

\usepackage[a4paper,margin=1in]{geometry}
\usepackage[T1]{fontenc}
\usepackage{lmodern}
\usepackage{microtype}
\usepackage{amsmath,amssymb,amsthm,mathtools}
\usepackage{booktabs,array}
\usepackage{enumitem}
\usepackage{graphicx}
\usepackage{xcolor}
\usepackage[colorlinks=true,linkcolor=blue!55!black,citecolor=blue!55!black,urlcolor=blue!55!black]{hyperref}

\newtheorem{theorem}{Theorem}
\newtheorem{remark}{Remark}
\newcommand{\cube}{\{0,1\}^n}
\newcommand{\E}{\mathbb{E}}
\newcommand{\wh}{\widehat}
\newcommand{\symdiff}{\mathbin{\triangle}}

\title{A Fourier-Positive Even-Exponent Subcase\\
of the Reverse Bernstein Problem on the Hypercube}
\author{Automatic Math Discovery Run \texttt{fa\_banach\_001}\\
Lane 8 (\texttt{agent\_lane\_08}), model GPT5.6}
\date{19 July 2026}

\begin{document}
\maketitle

\begin{center}
\fbox{\parbox{0.91\textwidth}{
\textbf{Status: candidate partial result, likely valid.}
The general scalar reverse Bernstein problem for arbitrary Fourier signs and arbitrary
$p\in(1,\infty)\setminus\{2\}$ remains open.  This packet proves the sharp estimate
for every even exponent when all Walsh coefficients can be made nonnegative by a
translation of the cube.
}}
\end{center}

\section{Source and original problem}

The source is P. Ivanisvili, R. Vershynin, and X. Xie,
\emph{Jackson's inequality on the hypercube}, arXiv:2410.19949v1
\cite{IVX}.  The problem appears in Section 1.3, page 3 of the arXiv PDF
(source file \texttt{source.tex}, line 314).  In the source normalization,
\[
  \Delta f(x)=\sum_{j=1}^n\bigl(f(x)-f(x^j)\bigr),
\]
where $x^j$ is obtained by flipping coordinate $j$.  Thus
$\Delta W_S=2|S|W_S$ for the Walsh character
\[
  W_S(x)=(-1)^{\sum_{j\in S}x_j}.
\]
The tail space $L^p_{\ge d}$ consists of the functions whose Walsh coefficients
vanish on all levels $|S|<d$.

The source states that, for $p\in(1,\infty)$, the reverse Bernstein estimate
\begin{equation}\label{eq:open}
  \|f\|_p \leq \frac{C_p}{d}\|\Delta f\|_p,
  \qquad f\in L^p_{\ge d},
\end{equation}
remains a major open problem.  The case $p=2$ is immediate from Parseval, with
$C_2=1/2$, so the substantive open range is $p\ne2$.

\begin{figure}[ht]
  \centering
  \includegraphics[width=\textwidth]{figures/open_problem_crop.png}
  \caption{Full-width crop of Section 1.3 on page 3 of the source paper.}
\end{figure}

\section{Definitions and scope of the partial result}

For $y\in\cube$, let $\tau_y f(x)=f(x\oplus y)$, where $\oplus$ denotes
coordinatewise addition modulo two.  Since
\[
  \wh{\tau_y f}(S)=W_S(y)\wh f(S),
\]
we make the following definition.

\medskip
\noindent
\textbf{Fourier-positive up to translation.}
A real-valued $f:\cube\to\mathbb R$ has this property if there is a
$y\in\cube$ such that
\[
  W_S(y)\wh f(S)\ge0
  \qquad\text{for every }S\subseteq[n].
\]
Equivalently, one translate of $f$ has nonnegative Walsh coefficients.

This is a genuine sign restriction.  It is preserved by multiplying $f$ by a
positive scalar and includes single Walsh characters after a translation and
choice of positive coefficient.  It does not include general tail functions.

\section{Idea of the proof}

For an even exponent $p=2m$, both $\|f\|_p^p$ and the dual pairing
$\E[f^{p-1}\Delta f]$ can be expanded as sums over $2m$ Walsh frequencies whose
symmetric difference is empty.  When the Walsh coefficients are nonnegative,
every summand is nonnegative.  Symmetry among the $2m$ frequencies lets us replace
the Laplacian weight on the last frequency by the average of all $2m$ weights.
Every active frequency has size at least $d$, so that average is at least $2d$
in the source normalization.  Holder's inequality then converts the pairing
estimate into the desired norm estimate.  A single level-$d$ character shows
that the constant $2d$ is sharp.

The positivity is precisely what prevents cancellations in the multilinear
Walsh expansion.  Removing it is the unresolved obstacle; the argument does not
claim progress for arbitrary Fourier signs.

\section{The partial theorem}

\begin{theorem}[Sharp Fourier-positive even-$p$ reverse Bernstein estimate]
\label{thm:main}
Let $n,d,m\in\mathbb N$ with $1\le d\le n$, and put $p=2m$.
Suppose that $f:\cube\to\mathbb R$ satisfies
\[
  \wh f(S)=0\quad\text{when }|S|<d
\]
and is Fourier-positive up to translation.  Then
\begin{equation}\label{eq:main}
  \|\Delta f\|_{p}\ge 2d\,\|f\|_{p}.
\end{equation}
The factor $2d$ is optimal for every even $p$.
\end{theorem}

\begin{proof}
Translations commute with $\Delta$ and preserve every $L^p$ norm.  We may
therefore translate $f$ and assume that
\[
  a_S:=\wh f(S)\ge0
  \qquad\text{for every }S\subseteq[n].
\]
The tail hypothesis says that $a_S=0$ whenever $|S|<d$.

Write $p=2m$.  Orthogonality of Walsh characters gives
\begin{equation}\label{eq:moment}
  \E f^{2m}
  =
  \sum_{S_1\symdiff\cdots\symdiff S_{2m}=\varnothing}
  \prod_{j=1}^{2m} a_{S_j}.
\end{equation}
Indeed, the expectation of $W_{S_1}\cdots W_{S_{2m}}$ is one exactly when
$S_1\symdiff\cdots\symdiff S_{2m}=\varnothing$, and is zero otherwise.

Since $\Delta W_S=2|S|W_S$, the same expansion yields
\begin{align}
  \E\bigl[f^{2m-1}\Delta f\bigr]
  &=
  \sum_{S_1\symdiff\cdots\symdiff S_{2m}=\varnothing}
  2|S_{2m}|\prod_{j=1}^{2m} a_{S_j}. \label{eq:weighted}
\end{align}
The indexing condition and the coefficient product are symmetric in
$S_1,\ldots,S_{2m}$.  Consequently, the value of the sum in
\eqref{eq:weighted} is unchanged if $|S_{2m}|$ is replaced by $|S_r|$ for
any $r$, and hence
\begin{align}
  \E\bigl[f^{2m-1}\Delta f\bigr]
  &=\frac{1}{2m}
  \sum_{S_1\symdiff\cdots\symdiff S_{2m}=\varnothing}
  2\Bigl(\sum_{j=1}^{2m}|S_j|\Bigr)
  \prod_{j=1}^{2m}a_{S_j}. \label{eq:symmetrized}
\end{align}
Every product in \eqref{eq:symmetrized} is nonnegative.  If it is nonzero,
then $a_{S_j}>0$ for every $j$, so the tail hypothesis gives $|S_j|\ge d$
for every $j$.  Therefore
\[
  \frac{2}{2m}\sum_{j=1}^{2m}|S_j|\ge 2d
\]
term by term.  Comparing \eqref{eq:symmetrized} with \eqref{eq:moment},
we obtain
\begin{equation}\label{eq:pairing}
  \E\bigl[f^{p-1}\Delta f\bigr]
  \ge 2d\,\E f^p
  =2d\,\|f\|_p^p,
\end{equation}
where the last equality uses that $p$ is even.

Holder's inequality and \eqref{eq:pairing} now imply
\[
  2d\,\|f\|_p^p
  \le \E\bigl[f^{p-1}\Delta f\bigr]
  \le \|f^{p-1}\|_{p/(p-1)}\|\Delta f\|_p
  =\|f\|_p^{p-1}\|\Delta f\|_p.
\]
If $f\ne0$, division by $\|f\|_p^{p-1}$ proves \eqref{eq:main}; the zero
case is trivial.

Finally, take $f=W_S$ with $|S|=d$.  Then $|f|=1$ and
$\Delta f=2d f$, so equality holds in \eqref{eq:main}.  This proves
sharpness.
\end{proof}

\begin{remark}[Alternative Laplacian normalization]
Some cited papers use $L=\Delta/2$, so that $LW_S=|S|W_S$.  In that
normalization Theorem \ref{thm:main} reads $\|Lf\|_p\ge d\|f\|_p$.
\end{remark}

\section{Adversarial verification}

\subsection*{Step check}

\begin{center}
\begin{tabular}{@{}>{\raggedright\arraybackslash}p{0.25\textwidth}>{\raggedright\arraybackslash}p{0.14\textwidth}>{\raggedright\arraybackslash}p{0.52\textwidth}@{}}
\toprule
Step & Verdict & Check \\
\midrule
Translation reduction & Valid & $\tau_y$ is a cube automorphism, commutes with $\Delta$, preserves the uniform measure, and changes coefficient signs by $W_S(y)$. \\
Walsh moment expansion & Valid & The product of characters is the character indexed by the symmetric difference; its mean is zero unless that index is empty. \\
Termwise positivity & Valid & This is exactly where nonnegative translated Fourier coefficients are used.  Zero coefficients cause no difficulty. \\
Symmetrization of weights & Valid & The constraint and coefficient product are invariant under all permutations of the $2m$ frequency indices. \\
Tail lower bound & Valid & A nonzero product forces all $2m$ participating coefficients to be nonzero, hence all frequency sizes are at least $d$. \\
Holder and sharpness & Valid & The dual exponent is $p/(p-1)$, and a single level-$d$ character gives equality. \\
\bottomrule
\end{tabular}
\end{center}

\subsection*{Finite checks}

The reusable script
\texttt{code/verify\_fourier\_positive\_even\_p.py} generated independent
exponential nonnegative coefficients for every $2\le n\le6$, every
$1\le d\le n$, and $p\in\{4,6,8\}$.  Ten trials per parameter triple gave
600 checks and no violation.  The smallest reported slack in
$\|\Delta f\|_p/\|f\|_p-2d$ was $-1.78\times10^{-15}$, numerical roundoff at
the equality case $d=n$.  The command was
\begin{verbatim}
conda run --no-capture-output -n sandbox python \
  code/verify_fourier_positive_even_p.py --max-n 6 --trials 10
\end{verbatim}
These checks are not used in the proof.

The separate script \texttt{code/finite\_tail\_inverse\_probe.py} performs a
nonconvex numerical search without the sign restriction.  For $n\le6$ and
$d$ near $n/2$, probes at $p=1.5$ and $p=4$ kept
$d\|\Delta^{-1}\|$ near $0.51$--$0.55$, compared with the exact Hilbert value
$1/2$.  This found no small counterexample to the full conjecture, but it
cannot certify an operator norm or prove the conjecture.

\subsection*{Verifier verdict}

\textbf{Likely valid partial result.}  The proof is finite, algebraic, and has
no external mathematical dependency beyond elementary Walsh orthogonality and
Holder's inequality.  Human review should focus on the permutation averaging
in \eqref{eq:symmetrized} and the factor two coming from the source paper's
definition of $\Delta$.

\section{Failed full-solution routes and remaining obstruction}

Three broader routes were tested before isolating Theorem \ref{thm:main}.
\begin{enumerate}[leftmargin=*,itemsep=0.4em]
  \item Integrating the best known pure-$L^p$ heat-smoothing estimate gives only
  a power $d^{-\alpha_p}$ with $\alpha_p<1$ for $p\ne2$, not $d^{-1}$.
  \item The nonlinear $L^p$ Poincare pairing controls mean-zero functions but
  does not see the tail depth $d$; inserting it loses the required factor.
  \item A tensor-product construction does not immediately yield a
  counterexample because the tilted mean of the block logarithmic Laplacian is
  positive.  Making that mean small is essentially another form of the open
  inequality.
\end{enumerate}

For arbitrary Fourier signs, the sums in \eqref{eq:moment} and
\eqref{eq:weighted} contain cancellations.  The weighted sum need not dominate
the unweighted moment term by term.  No argument in this packet controls those
cancellations, and no claim is made for non-even $p$.

\section{Novelty and literature-status check}

The bounded check was performed on 19 July 2026 after the direct proof attempt.
It covered the run's four lightweight indexes, the 35,297-paper local
source-probed corpus, exact-phrase searches for ``reverse Bernstein'' and
``tail space'', searches for the heat-smoothing conjecture, and searches for
the combinations ``nonnegative Fourier coefficients'', ``Fourier-positive'',
``Hamming cube'', and ``Laplacian''.  The close papers found were:
\begin{itemize}[leftmargin=*]
  \item Eskenazis--Ivanisvili \cite{EI}, which proves sublinear-exponent bounds
  in general and a sharp-order result for narrow spectra;
  \item later work on complex hypercontractivity and the heat-smoothing
  exponent, including arXiv:1907.11359;
  \item Domelevo et al. \cite{DDF}, which proves linear heat decay with mixed
  $L^2$ and neighboring-$L^p$ norms, not the pure estimate \eqref{eq:open}.
\end{itemize}
No exact match for Theorem \ref{thm:main} was found.  This bounded search does
not establish novelty, so novelty confidence is \emph{moderate}.  The theorem
should be treated as a candidate new partial result until a specialist checks
the Fourier-analytic literature.

\section{Limitations and human-review recommendation}

Theorem \ref{thm:main} settles an infinite, sharp subcase of the scalar
question, but it does not settle the general problem.  Its two restrictions
are substantial:
\begin{itemize}[leftmargin=*]
  \item $p$ must be an even positive integer;
  \item after a cube translation, all Walsh coefficients must be nonnegative.
\end{itemize}
The recommended human action is to verify the four-line multilinear expansion
and then assess whether the positivity mechanism is known under a different
name (for example, positive-definite functions on the Boolean group).  A
secondary research direction is to replace coefficientwise positivity by a
weaker $2m$-fold signed-energy condition that still makes
\eqref{eq:symmetrized} coercive.

\enlargethispage{4\baselineskip}
\small
\begin{thebibliography}{9}

\bibitem{IVX}
P. Ivanisvili, R. Vershynin, and X. Xie,
\emph{Jackson's inequality on the hypercube},
arXiv:2410.19949v1, 2024.

\bibitem{EI}
A. Eskenazis and P. Ivanisvili,
\emph{Polynomial inequalities on the Hamming cube},
Probability Theory and Related Fields 178 (2020), 235--287;
arXiv:1902.02406.

\bibitem{DDF}
K. Domelevo, P. Durcik, V. Fragkiadaki, O. Klein,
D. Oliveira e Silva, L. Slavikova, and B. Wrobel,
\emph{Dimension-free estimates for low degree functions on the Hamming cube},
Studia Mathematica 280 (2025), no. 2, 161--176;
arXiv:2401.07699.

\end{thebibliography}

\end{document}
